\documentclass[main.tex]{subfiles}
\begin{document}
\ifSubfilesClassLoaded{\begin{center}\small\itshape Part of the GnssProcessor Mathematical Reference \textbullet\ \today\end{center}}{}
\section{Real-time SSR corrections}\label{sec:ssr}
% ============================================================
%  Real-time corrections: the SSR-to-observation mapping
% ============================================================

The real-time engine consumes network corrections in \emph{State Space
Representation} (SSR): per-satellite, per-signal refinements to the
broadcast orbit, clock and hardware biases. They arrive as RTCM~3.x
messages (legacy 1057--1068 and 1240--1270) or in the IGS-SSR
\texttt{4076} repackaging \cite{rtcm10403,igsssr}. Each SSR quantity
refines one term already present in
\dcref{eq:obs-code}{the undifferenced code equation} and
\dcref{eq:obs-phase}{the undifferenced phase equation} of the
Foundations document, and none introduces a new state. SSR is what
distinguishes the precise real-time mode from broadcast-only
operation: with SSR applied, $\sLight\,\mathrm{d}t^\sat$ and $\satpos$
are centimetre-grade, and the satellite biases $B^\sat_\sigl,
b^\sat_\sigl$ are \emph{known} instead of absorbed into the ambiguity.
This document describes the user-side application of the corrections. The
same store receives a third-party IGS-RTS stream and the products of the
engine's own network estimator (\dcref{sec:net}{the Network mode
document}), so the application mathematics below does not depend on their
source.

Decoded records are cached per satellite in
\texttt{products/SsrCorrectionStore}: a process-wide singleton for the
live stream, or a caller-owned instance for the closed-loop user of the
network (\Cref{sec:ssr-netuser}). The orbit is applied in
\texttt{model/ModelGeometry}, the clock in
\texttt{corrections/GnssCorrections} and the biases in
\texttt{measurement/ObservationEquation}. Each queries the store and falls
back to broadcast-only when it holds nothing usable. In post-processing the store is
empty, so every equation below reduces to the broadcast-only observation
model. The streaming and decoding layer (RTCM frame parsing, the
per-satellite mutex-guarded cache, the real-time push loop) is described in
\dcref{sec:rt}{the Real-time engine document}. This document covers only
the store-driven correction mathematics, which does not depend on how the
records arrived. There is no \texttt{enable\_ssr} flag: SSR is applied
whenever the store holds a usable record.

\subsection{Orbit correction: refines $\geomrange^\sat_\rx$}

The orbit message carries a delta in the satellite's
radial, along-track and cross-track (RAC) frame, referenced to the
broadcast ephemeris of issue $\mathrm{IODE}$. The frame is built from
the broadcast position $\satpos_{\mathrm{brdc}}$ and velocity
$\satvel_{\mathrm{brdc}}$ (\texttt{products/SsrCorrectionStore.cpp}):
\begin{equation}
\racA = \frac{\satvel_{\mathrm{brdc}}}{\lVert\satvel_{\mathrm{brdc}}\rVert},
\quad
\racC = \frac{\satpos_{\mathrm{brdc}}\times\satvel_{\mathrm{brdc}}}
             {\lVert\satpos_{\mathrm{brdc}}\times\satvel_{\mathrm{brdc}}\rVert},
\quad
\racR = \racA\times\racC .
\label{eq:ssr-rac-frame}
\end{equation}
The RAC delta is propagated to the query epoch with its rate term and
rotated into ECEF. With $\delta t = t - t_{\mathrm{SSR}}$ the age of the
correction relative to the content epoch of the orbit record,
$\bm{\delta}_{\mathrm{rac}} = (\delta_r,\delta_a,\delta_c)^\Tr$
its components and $\dot{\bm{\delta}}_{\mathrm{rac}}$ their rate,
\begin{empheq}[box=\fbox]{equation}
\satpos
\;=\;
\satpos_{\mathrm{brdc}}
\;-\;
\big[\,\racR\;\;\racA\;\;\racC\,\big]
\big(\bm{\delta}_{\mathrm{rac}} + \dot{\bm{\delta}}_{\mathrm{rac}}\,\delta t\big).
\label{eq:ssr-orbit}
\end{empheq}
The subtraction follows the RTCM sign convention
$\bm{\delta} = \text{broadcast} - \text{precise}$
(\texttt{products/SsrCorrectionStore.cpp}). If the broadcast velocity or
the position-velocity cross product is degenerate, the frame is
undefined and the broadcast position is returned unchanged. The
corrected $\satpos$ enters $\geomrange^\sat_\rx = \lVert\satpos -
\rxpos_\rx\rVert$ and the line-of-sight $\los^\sat_{\hat{\rx}}$ of the
Foundations document unchanged. The orbit correction therefore acts on no
observation equation directly, only on the geometry
(\dcref{sec:obs}{the observation-geometry section}) that feeds it.

The delta is applied only when its $\mathrm{IODE}$ matches that of the
broadcast ephemeris from which the baseline orbit was evaluated
(\Cref{eq:ssr-iode-gate}). Over a product orbit the gate compares with the
IODE of the orbit node nearest in time; SP3 nodes carry none, so the gate
passes. On a live ephemeris cutover, the baseline
position is re-referenced to the ephemeris \emph{version} against which
the SSR delta was issued, so the satellite does not drop to broadcast
grade (\texttt{model/ModelGeometry.cpp}, via
\texttt{LiveEphemerisStore::computeSatellitePV} with the SSR IODE).
The clock is queried against the same broadcast ephemeris, so the
broadcast clock's \emph{radial} component cancels consistently against
\Cref{eq:ssr-clock}. The decoded RAC scale factors are those of
RTCM~10403.3: the radial $\delta_r$ is a $22$-bit two's-complement field at
$0.1$\,mm, along-track and cross-track are at $0.4$\,mm, and the rates are
at $1$ and $4\,\mu$m/s respectively (\texttt{rtcm/RtcmSsrDecoder.cpp}).

The orbit record also carries a \emph{reference point} for the corrected
$\satpos$. Legacy RTCM SSR corrects to the satellite antenna phase centre
(APC); IGS-SSR~\texttt{4076} and the network product correct to the centre
of mass (CoM) (\texttt{rtcm/RtcmSsrDecoder.cpp}). Under
\texttt{realtime.honor\_ssr\_reference\_point} (default OFF,
byte-identical) the geometry records whether the applied overlay was
APC-referenced (\texttt{ModelGeometry::ssrOrbitApcApplied}). \texttt{PppModel} then omits its own
satellite-PCO term on that geometry, so the phase-centre offset already
contained in the orbit correction is not counted twice.

\subsection{Clock correction: refines $-\sLight\,\mathrm{d}t^\sat$}

The clock message carries a quadratic polynomial in the time since the
reference epoch, optionally augmented by an additive high-rate delta
$\delta C_{\mathrm{hr}}$. The total satellite-clock correction, in metres, is
\begin{equation}
\Delta C^\sat(t)
\;=\;
C_0 + C_1\,\delta t + C_2\,\delta t^{2} + \delta C_{\mathrm{hr}},
\qquad \delta t = t - t_{\mathrm{SSR}},
\label{eq:ssr-clock-poly}
\end{equation}
converted to seconds (\texttt{SsrCorrectionStore::applyClockToBias}) and
\emph{added} to the baseline clock
(\texttt{corrections/GnssCorrections.cpp}):
\begin{empheq}[box=\fbox]{equation}
\sLight\,\mathrm{d}t^\sat
\;=\;
\sLight\,\mathrm{d}t^\sat_{\mathrm{base}}
\;+\;
\Delta C^\sat(t),
\label{eq:ssr-clock}
\end{empheq}
where $\mathrm{d}t^\sat_{\mathrm{base}}$ is the precise CLK sample
(post-processing) or the broadcast clock polynomial (live mode,
\texttt{products/LiveEphemerisStore}). The term enters both
undifferenced observation equations through $-\sLight\,\mathrm{d}t^\sat$,
identically on code and phase. The high-rate clock is added to the
polynomial and does not replace it
(\texttt{corrections/GnssCorrections.cpp}); it is subject to its own
freshness test in the store. The $C_2$ field is a
$27$-bit two's-complement quantity at $2\times10^{-8}$\,m/s$^2$. The time
base of the polynomial is the record's content epoch (the provider's
reference epoch), not its arrival time.

\subsection{IOD gate at correction-set boundaries}\label{sec:ssr-iod}

Orbit and clock corrections are deltas against a specific
broadcast ephemeris, so they are applied only when the SSR record's
$\mathrm{IODE}$ matches that of the broadcast ephemeris from which the
baseline orbit/clock was evaluated:
\begin{equation}
\text{apply}
\;\Longleftrightarrow\;
\big(\mathrm{IODE}_{\mathrm{brdc}} < 0\big)
\;\vee\;
\big(\mathrm{IODE}_{\mathrm{SSR}} < 0\big)
\;\vee\;
\big(\mathrm{IODE}_{\mathrm{brdc}} = \mathrm{IODE}_{\mathrm{SSR}}\big),
\label{eq:ssr-iode-gate}
\end{equation}
The gate is permissive when either side is unknown
(\texttt{SsrCorrectionStore::iodeMatches},
\texttt{products/SsrCorrectionStore.hpp}). A combined orbit+clock
record (RTCM 1060/1066, IGS-SSR sub-type~3) takes precedence over separate
messages because it guarantees a single shared $\mathrm{IODE}$
(\texttt{products/SsrCorrectionStore.cpp}). The two overlays gate on
separate anchors: the orbit delta on the record's orbit
\texttt{iode} (\texttt{model/ModelGeometry.cpp}), the clock polynomial on
a distinct \texttt{clock\_iode} field
(\texttt{corrections/GnssCorrections.cpp}). For a standard stream,
\texttt{clock\_iode} inherits the orbit anchor in \texttt{getOrbitClock},
so the two coincide and \eqref{eq:ssr-iode-gate} is a single gate. Only a
network clock that stamps its own clock IODE (the \texttt{brdc\_rts} clock
mode) makes them differ. At a correction-set boundary (the provider switches to a new
ephemeris reference, hence a new $\mathrm{IODE}$), the gate
\eqref{eq:ssr-iode-gate} declines the mismatched record until the matching
broadcast ephemeris is available. This avoids a discontinuous jump in
$\satpos$ or $\sLight\,\mathrm{d}t^\sat$.

The two sides of \eqref{eq:ssr-iode-gate} are compared in the IGS-SSR
convention. A \texttt{4076} message carries an 8-bit $\mathrm{IODE}$ field
for every constellation. A legacy Galileo message carries the full 10-bit
IODnav, and the decoder keeps its low 8 bits
(\texttt{normalizeSsrIode} in \texttt{rtcm/RtcmSsrDecoder.cpp}). The
broadcast side converts its record in the same way
(\texttt{LiveEphemerisStore::ssrIode}): the GPS and QZSS IODE as broadcast,
the low 8 bits of the Galileo IODnav, and
$\lfloor t_{oe}/720\rfloor \bmod 240$ for BeiDou, with $t_{oe}$ the
seconds of week of the ephemeris reference time in BDT. The BeiDou
broadcast AODE is a different quantity and is not used.

Records older than \texttt{max\_age\_s}
(\texttt{processing.ssr\_max\_age\_s}, default $30$\,s) are treated as
absent. The age is that of the record's receive time relative to the
query epoch, and the test is symmetric: a record dated more than
the same limit ahead of the observation (a desynchronised
live stream) is also rejected. The stale-hold rule extends this
binary clock gate to three stages
(\texttt{SsrCorrectionStore::classifyClockAge}). With
\texttt{processing.ssr\_clock\_stale\_hold\_factor}~$f>1$, a clock aged in
$(\texttt{max\_age\_s},\,f\cdot\texttt{max\_age\_s}]$ is \emph{StaleHeld}.
It is still applied, extrapolated over the full age: the code evaluates
the whole polynomial \eqref{eq:ssr-clock-poly}, of which
$C_0 + C_1\,\delta t$ is the dominant part. The variance of the consuming
rows grows by $q\,\delta t^{3}/3$, the zero-order-hold variance of a clock
that wanders with density $q$ (\texttt{network.ura\_staleness\_q2}, which
defaults to \texttt{network.clock\_model\_q2}; $1.5\times10^{-7}$\,m$^2$/s$^3$
is used when it is zero;
\texttt{measurement/MeasurementModelVariance.cpp}). A record beyond the
hold window is \emph{Rejected}. The default $f=0$ reduces the rule
to the plain fresh-or-absent gate (byte-identical).

\paragraph{Orbit-coherence echo.}
A network clock in the \texttt{brdc\_rts} mode is estimated against one
orbit update, and the radial error of the broadcast orbit is absorbed one for
one into that clock. A user that applies a different orbit update receives a
range bias of decimetres to metres that residual tests do not detect. The
network therefore echoes the identity of the orbit it used (its IOD-SSR and
content epoch) on each clock record, and the user compares it with the orbit it
applies (\texttt{SsrCorrectionStore::classifyOrbitEcho}). The verdict is one
of: no echo stamped, match, no applied orbit, IOD-SSR mismatch, or
content-epoch skew beyond a budget. Every verdict other than a match is
counted per satellite. The clock is withheld only when
\texttt{processing.network\_consumer.orbit\_clock\_max\_skew\_s} is positive
(default $0$: count only), so a stream that stamps no echo is unaffected.

\subsection{Satellite biases: refine $-B^\sat_\sigl$ and $-b^\sat_\sigl$}

SSR delivers \emph{observable-specific biases} (OSB): one value per
(satellite, RINEX-3 signal) pair, in metres. They are applied as a
known additive correction to the modelled observable, with no
state and no Jacobian column
(the \texttt{SsrBias} term built by \texttt{addSsrBias} in
\texttt{measurement/ObservationEquation.cpp}). The code OSB refines
$B^\sat_\sigl$ in the undifferenced code equation:
\begin{empheq}[box=\fbox]{equation}
\code^\sat_{\rx,\sigl}
=
\geomrange^\sat_\rx + \sLight(\mathrm{d}t_\rx - \mathrm{d}t^\sat)
+ \sLight\,\delta t^{(g)}_\rx + T^\sat_\rx + \mu_\sigl I^\sat_{\rx,1}
+ B_{\rx,\sigl} \underbrace{-\,B^\sat_\sigl}_{\text{SSR OSB}}
+ \varepsilon^\code,
\label{eq:ssr-code-bias}
\end{empheq}
and the phase OSB refines $b^\sat_\sigl$ in the undifferenced phase
equation:
\begin{empheq}[box=\fbox]{equation}
\phaseo^\sat_{\rx,\sigl}
=
\geomrange^\sat_\rx + \sLight(\mathrm{d}t_\rx - \mathrm{d}t^\sat)
+ \sLight\,\delta t^{(g)}_\rx + T^\sat_\rx - \mu_\sigl I^\sat_{\rx,1}
+ \lambda_\sigl N^\sat_{\rx,\sigl}
+ b_{\rx,\sigl} \underbrace{-\,b^\sat_\sigl}_{\text{SSR OSB}}
+ \varepsilon^\phaseo .
\label{eq:ssr-phase-bias}
\end{empheq}
The stored value is the satellite bias as it is present in the
observation (the underbraced term), so it is added to the clean model
prediction with sign $+1$
(\texttt{kSatBiasSign} in \texttt{measurement/ObservationEquation.hpp}),
the constant that is also the partial of the network estimator's satellite
code-bias and phase-bias states. Code and phase share it. The sign relating
the value a provider sends to the stored value is fixed in the bias bodies of
\texttt{rtcm/RtcmSsrDecoder.cpp}. It can be confirmed only empirically
against a provider's stream: float ambiguity fractions that cluster near one
half instead of zero indicate an inverted constant. The value enters the
predicted observable as the \texttt{SsrBias} term of the observation
equation, summed with the receiver clock, troposphere, ionosphere,
estimated code and phase biases and the ambiguity.
Biases are queried per \emph{raw} signal only: an ionosphere-free
composite key such as \texttt{C1W\_C5X} returns zero (code) or
\texttt{nullopt} (phase) and is skipped, consistent with the uncombined
parameterisation. The code OSB is a $14$-bit field at $0.01$\,m, the
phase OSB a $20$-bit field at $0.1$\,mm
(\texttt{rtcm/RtcmSsrDecoder.cpp}).

The 5-bit signal identifier of each bias is translated through the SSR
table (\texttt{RtcmSignalMap::ssrToRinexSuffix}), which differs from the MSM
table that numbers the cells of an observation message. GPS L1 C/A is SSR
identifier $0$ and L2 Z-tracking is $11$, against MSM identifiers $2$ and
$10$; only a few identifiers coincide (Galileo E1 C is $2$ in both)
(\texttt{rtcm/RtcmSignalMap.hpp}). A bias whose identifier has no entry for
its constellation is dropped by the decoder, and a bias read through the
wrong table would be applied to the rows of another signal.

By default the application path has no age gate: the bias
is loaded whenever the store carries the exact RINEX signal key (or, for
code only and with the optional band fallback
\texttt{processing.ssr\_bias\_band\_fallback} armed, a same-band
sibling key; phase never falls back, to preserve the integer property). A
lookup that misses the exact key while a same-band sibling exists is counted
per satellite, signal and domain, whether or not the fallback resolves it.
The optional staleness gate (\texttt{processing.enable\_ssr\_bias\_max\_age},
default OFF) measures the age of the bias record against the observation
epoch. Its limit is \texttt{ssr\_bias\_stale\_hold\_factor}~$\times$
\texttt{processing.ssr\_bias\_max\_age\_s}; the factor defaults to 2 and the
maximum age defaults to the clock maximum age, hence $30$\,s. Below that limit the
last value stays applied at full weight across staleness gaps, so a
provider update interval near the maximum age produces no model step; a
bias older than the limit is dropped (\texttt{measurement/ObservationEquation.cpp}).
The derived configuration of the closed-loop user and the standalone
\texttt{network\_consumer} profile arm the gate under bias consumption unless the
configuration sets the key explicitly.
The consequence of the application path for the state vector, avoiding a
redundant estimated bias on the same row, is the subject of the next
subsection.

\subsection{Satellite-bias state gate and SSR code biases}\label{sec:ssr-satbias-gate}

The uncombined model can also \emph{estimate} a per-(satellite, signal)
code-bias state (\texttt{estimation.estimate\_sat\_biases}; see
\dcref{sec:state}{the state-vector section}). The estimated state and a
known SSR code OSB both enter the code row with unit coefficient, so
allowing both for the same (\sat, \sigl) is redundant: the stream already
fixes the satellite code bias, and a free state collinear with it is
nearly unobservable and absorbs noise. State allocation therefore consults
both correction sources before creating a sat-bias state
(\texttt{MeasurementModel::ensureRequiredStates} in
\texttt{measurement/MeasurementModelStates.cpp}):
\begin{equation}
\text{create sat-bias}(\sat,\sigl)
\;\Longleftrightarrow\;
\neg\,\mathrm{hasBIA}(\sat,\sigl)
\;\wedge\;
\neg\,\mathrm{hasSsrCB}(\sat,\sigl)
\;\wedge\;
\mathrm{pivotCovers}(\sat,\sigl),
\label{eq:ssr-satbias-gate}
\end{equation}
where $\mathrm{hasBIA}$ tests the static \texttt{.BIA} product
(\texttt{biaCoversSignal}, from \texttt{PppModel::hasBiasEntry}),
$\mathrm{hasSsrCB}$ tests the live stream
(\texttt{SsrCorrectionStore::hasCodeBias}, surfaced as
\texttt{ssrCoversCodeSignal} in \texttt{measurement/MeasurementModelRow.cpp})
and $\mathrm{pivotCovers}$ is the band gate of between-satellite
differencing (\texttt{signalCoveredByPivot}): only signals that the pivot
satellite tracks are observable in a single difference, and a state created
for another signal would be seeded against an unconverged state vector, so
its creation waits until the signal is usable. On a GLONASS FDMA signal that carries a
per-slot inter-frequency-bias state, no sat-bias state is created either,
because the two would be collinear. The gate rests on two design points:
\begin{itemize}[leftmargin=*]
  \item \textbf{Presence, not magnitude.} \texttt{hasCodeBias} tests
        whether the key is \emph{present} in the per-signal map, which
        differs from \texttt{getCodeBias}~$\neq 0$: a provider's
        present-but-zero entry still defines that signal's datum, so the
        estimator must not re-estimate it. The two queries agree by
        construction, including the optional same-band fallback, so a
        covered signal is never both estimated and SSR-corrected on one
        row. With the staleness gate armed, the presence test uses the
        same age limit as the application path; otherwise it ignores age.
  \item \textbf{Per-signal, not global.} IGS streams are partial (for
        example, Galileo E1 covered and E5a/E5b not). The gate keys on the
        exact (\sat, \sigl) pair. A signal the stream does not carry keeps
        its estimated state, its only available code correction, while a
        covered signal on the same satellite is left to the OSB. A strict
        per-epoch prune (\texttt{pruneSatBiasStatesExcept}) drops the state
        when SSR coverage arrives mid-arc and re-creates it (warm-restarted
        from the stashed cache) if coverage later disappears.
\end{itemize}
In post-processing the live store is empty, so
$\mathrm{hasSsrCB}$ is always false and the gate reduces to the
static-\texttt{.BIA} test.

\paragraph{Exemption of the network estimator.} When the engine is the
network estimator (\texttt{ProcessingMode::PPP\_RTK\_Network}), the roles
are reversed: it \emph{estimates} the satellite code and phase biases as
states and emits them as a product (\dcref{sec:net}{the Network mode
document}). Importing SSR biases there would count them twice and make
those states unobservable. Both the application
path \eqref{eq:ssr-code-bias}, \eqref{eq:ssr-phase-bias}
(\texttt{addSsrBias} adds no term) and the $\mathrm{hasSsrCB}$ branch of the
state gate \eqref{eq:ssr-satbias-gate} (\texttt{ssrCoversCodeSignal} returns
false) are therefore disabled for the network, whose own admission owns the
satellite code-bias states. The network's emitter
(\texttt{products/NetworkProductEmitter.cpp}, called from
\texttt{network/emit/EmitStore.cpp}) pushes the
estimated clocks and code and phase biases into a separate store instance,
which closes the loop against a user at a test station
(\dcref{sec:net-closedloop}{the closed-loop section of the Network mode
document}).

\paragraph{Integer recoverability via consistent phase biases.}
Removing $b^\sat_\sigl$ restores the integer property of
$N^\sat_{\rx,\sigl}$ stated in the phase observation equation of the
Foundations document. With the satellite phase bias known, the remaining
non-integer part of the ambiguity is the receiver phase bias
$b_{\rx,\sigl}$ alone, which cancels under between-satellite
single differencing $\bsd{\,\cdot\,}$
(\dcref{sec:diff}{the differencing section}). The result is an
integer single-difference ambiguity admissible to AR. This defines a
\emph{consistent} phase-bias product in operational terms: the provider
must generate $b^\sat_\sigl$ on the same network datum whose integer
property the rover then uses. The phase-bias message therefore carries
three metadata fields per signal that AR depends on
(\texttt{rtcm/RtcmTypes.hpp}): an integer-preservation indicator, a
two-bit wide-lane integer indicator, and a four-bit \emph{discontinuity
counter}. The message header carries two further one-bit indicators, the
dispersive-bias and the Melbourne-W\"ubbena consistency indicators. They are
read once, after the solution identifier and before the satellite count, and
copied onto every record of the message (\texttt{decodeSsrHeader} in
\texttt{rtcm/RtcmSsrDecoder.cpp}); they are not repeated in the satellite
block. This holds for the legacy messages \texttt{1265} to \texttt{1270} and
for \texttt{4076}. A decoder that reads them per satellite shifts the count
and every block after the first. The satellite block also holds a yaw angle
($9$ bits, $180/256$ degrees per unit) and a yaw rate ($8$ bits, $180/8192$
degrees per second per unit), which are stored and not used.

\subsection{Phase-bias steps and the user's ambiguities}\label{sec:ssr-pbstep}

A phase bias that changes between correction sets moves the term that the
model adds to every phase row of the satellite. Let $\tilde{b}^\sat_\sigl$ be
the applied bias (zero when none is applied, which includes a bias dropped by
the staleness gate) and $A^\sat_\sigl = \lambda_\sigl N^\sat_{\rx,\sigl}$ the
float ambiguity state in metres. The predicted phase carries
$\tilde{b}^\sat_\sigl + A^\sat_\sigl$, both with coefficient $+1$. A change
$\Delta = \tilde{b}_k - \tilde{b}_{k-1}$ between consecutive epochs carries no
information about the position, and the phase model stays consistent only if
the ambiguity absorbs it exactly:
\begin{equation}
A^\sat_\sigl \;\leftarrow\; A^\sat_\sigl - \Delta,
\qquad \text{covariance unchanged.}
\label{eq:ssr-pb-translate}
\end{equation}

A step is signalled in two ways. An incremented \emph{discontinuity counter}
(\texttt{StateManager::noteSsrDiscontinuity}) declares that the provider's
bias jumped, which breaks the integer connection of that arc. Under the
staleness gate, a bias that changes between applied and dropped (a freshness
flip) steps the model term without any counter. A step can also occur with no
declaration. The user reacts in \texttt{RealTimeEngine::resetOnPhaseBiasDisc},
in this order.
\begin{enumerate}[leftmargin=*, itemsep=0pt]
\item A bias whose store entry carries a nonzero hold generation (the
      network's phase-bias hold) is claimed first by \texttt{reactToPbHold\_}.
      Its answer, including the presence edges of a held bias, is in
      \dcref{sec:net-wire-consumer}{the user's answer in the event-response
      section}.
\item With \texttt{processing.ssr\_phase\_bias\_step\_translate} (default OFF)
      every step with $|\Delta|$ above
      \texttt{processing.ssr\_phase\_bias\_step\_translate\_eps\_m} (default
      $5$\,mm) is translated by \eqref{eq:ssr-pb-translate}, declared or not
      (\texttt{PbStepTranslateLedger}). Inside the band the filter absorbs the
      move as noise and nothing is written.
\item Otherwise, with \texttt{processing.ssr\_phase\_bias\_disc\_carry}
      (default OFF), a declared step is translated by
      \eqref{eq:ssr-pb-translate} (\texttt{PbDiscCarryLedger}). A counter
      change and a freshness flip on one key in one epoch are one event and
      one translation.
\item Otherwise the ambiguity is reset: the state is re-seeded from code at
      the configured ambiguity $\sigma$ (\texttt{StateManager::resetAmbiguity}).
\end{enumerate}
Both translations record, for every (satellite, signal), the value the model
applied at the previous update, read as the model reads it (age-gated when the
staleness gate is armed), so $\Delta$ is the step the phase row actually took.
The declared-step carry reads an unapplied bias as zero, so a freshness flip
is translated by the whole bias. The every-step translation needs an applied
value on both sides and sends a declared flip to the reset. Both fall back to
the reset, with each reason counted, when the key has no recorded value (a
first sighting, or the first event after a restore), when the ambiguity basis
is narrow-lane and wide-lane (the shift would need the basis transform), when
the user is between-satellite differenced (its states are differences) or
when the ambiguity key is not in the state. The every-step translation counts
an undeclared step in these classes without resetting. Each event takes
exactly one counted outcome, and the run-end counters close. The record of
applied values is not part of the filter-state snapshot, so the first
declared step per key after a warm restore takes the reset
(\dcref{sec:rt-snapshot}{the filter-state snapshot section}).

Both settings are off in the reference configurations
(\texttt{config/golden/golden\_field\_refound.json} and
\texttt{config/golden/golden\_sim\_refound.json}) and carry the ledger
verdict KEEP\_OFF. The declared-step carry is exact and does not move the
emitted product, but failed two of five bars in its re-grade
(\texttt{docs/refound/CC-arm\_pb\_disc\_carry\_rf0hv\_2026-09-15.md}). The
every-step translation was superseded by the network's phase-bias hold,
which changes what the product declares
(\texttt{docs/refound/PBH\_phase\_bias\_hold\_2026-09-19.md};
\dcref{sec:net-wire-pb}{the emitted phase bias}). The reset is the phase-bias
counterpart of the orbit and clock IOD gate \eqref{eq:ssr-iode-gate}: both
protect state continuity across correction-set boundaries, one by declining a
mismatched delta, the other by resetting the dependent ambiguity. A change of
the network's solution-generation id, and a re-reference of its code-bias or
clock datum, are answered by resets and re-inflations described in
\dcref{sec:rt-epoch}{the per-epoch update of the Real-time engine document}.

\subsection{GLONASS}\label{sec:ssr-glonass}

GLONASS SSR is decoded and stored in the same way as the CDMA
constellations. The legacy block (orbit 1063, clock 1064, code bias
1065, combined 1066, URA 1067, high-rate clock 1068, phase bias 1266)
and the IGS-SSR \texttt{4076} subtypes in the GLONASS
message-indicator band ($\mathrm{IM}=41$--$47$) both map onto the
same orbit, clock and bias structures
(\texttt{rtcm/RtcmSsrDecoder.cpp}). The orbit, clock and bias
equations above apply unchanged; nothing in
\Cref{eq:ssr-orbit,eq:ssr-clock,eq:ssr-code-bias,eq:ssr-phase-bias} is
constellation-specific.

The GLONASS-specific difference is the time tag of the legacy messages;
the correction mathematics is shared. Legacy GLONASS SSR (\texttt{1063} to
\texttt{1068} and \texttt{1266}) carries its content epoch as a 17-bit
time-of-day field (DF389, seconds since the start of the
GLONASS\,UTC(SU)$+3$\,h day) instead of a 20-bit GPS time of week, and its
satellite identifiers have 5 bits. IGS-SSR~\texttt{4076} uses one layout for
every constellation (6-bit satellite identifier, 8-bit IOD, 20-bit
GPS-aligned time of week), so the conversion applies to the legacy messages
only. The decoder converts the legacy field to the GPS time scale through
\texttt{TimeUtils::glonassTodToGpsTime}
(\texttt{rtcm/RtcmSsrDecoder.cpp}), resolving the day from the parser's
last known GPS stream time. Without that conversion the SSR epoch would
be $\{0,0\}$; the store would then drop every GLONASS record as stale or
IOD-mismatched (the \texttt{max\_age\_s} gate that accompanies
\eqref{eq:ssr-iode-gate}), and the satellite would fall back to broadcast
grade without any error. After conversion, the per-record content
epoch drives $\delta t = t - t_{\mathrm{SSR}}$ in
\Cref{eq:ssr-orbit,eq:ssr-clock} as for GPS. Whether the
resulting GLONASS observable then \emph{contributes} to a
between-satellite-differenced solution is a separate estimator-side
option (the FDMA per-channel receiver delay does not cancel in SD); see
\dcref{sec:diff}{the differencing section}.

\subsection{Ionosphere: $\mu_\sigl I^\sat_{\rx,1}$ off the stream, via the slant channel}

The SSR standard defines a global VTEC spherical-harmonic message
(IGS-SSR IM\,$=201$) and regional STEC products, but this engine decodes
neither. The IGS-SSR classifier rejects IM\,$=201$ as unsupported
(\texttt{rtcm/RtcmSsrDecoder.cpp}), so nothing on an RTCM stream writes the
slant-iono state $I^\sat_{\rx,1}$. For a user of an RTCM stream, the
ionospheric term $\mu_\sigl I^\sat_{\rx,1}$ is therefore informed only by
(i) the per-satellite iono \emph{state} estimated by the filter and
(ii) the optional GIM/IONEX deterministic prior
(\texttt{atmosphere/GlobalIonosphereMap}, an independent VTEC$\to$STEC
mapping unrelated to the SSR stream).

The store has a dedicated undifferenced slant-iono channel
(\texttt{SsrSlantIono}, via \texttt{SsrCorrectionStore::pushSlantIono} /
\texttt{getSlantIono}) that closes the atmosphere loop outside the RTCM
stream. When the network estimates a spherical-harmonic VTEC field, its emit
seam (\texttt{network/emit/EmitStore}) evaluates the field
(\dcref{sec:atmos-shvtec}{the co-estimated VTEC field}) at the user's pierce
points and pushes a per-LOS slant delay and $1\sigma$. The user
(\texttt{RealTimeEngine::applyNetworkIonoPrior}, run after the measurement
update and before AR) folds each fresh record into an ionosphere-weighted
(BLUP) pseudo-observation on the matching slant-iono state
(\texttt{KalmanFilter::applyPseudoObservation}), with innovation
$I^{\mathrm{net}}_i - \hat{I}_i$ and $\jacH$ selecting the states. This is the
ionosphere-weighted constraint that makes single-epoch AR feasible. The
fold is gated by \texttt{network.closed\_loop.consume\_network\_iono}
(default OFF: no rows are built and the update is byte-identical). It admits a
value only when the network refreshed it in the current epoch
(\texttt{max\_age}~$=\tfrac12$ the epoch spacing, or 30\,s when the spacing
is unknown), so a satellite the network did not re-supply floats instead
of inheriting a stale delay.

The reference configurations estimate no spherical-harmonic field (the slant
ionosphere is free per line of sight), so the supply is gated on a model that
is never built and the channel stays empty, although the reference
configurations carry the setting as true by lineage. The ledger verdict of the
setting is KEEP\_OFF (\texttt{docs/CONFIG\_promotion\_ledger.json};
\texttt{docs/refound/AB1\_sh\_consumer\_product\_2026-09-20.md}). The channel
has no record in the product file either, so the offline replay and the
file-tail user cannot receive it; the setting is inert there and draws a
warning.

Three properties of the channel follow from the field and are not chosen
by the user. The supply covers GPS and Galileo only: GLONASS is
excluded from the SH field, so a user's GLONASS slant states float.
When a second shell is armed, the pushed value is the \emph{sum} over
shells, each evaluated at its own pierce geometry, so the user does not
depend on the layer model. The per-LOS $\sigma$ is the field's
representativeness floor (\dcref{eq:atmos-sh-sigma}{the SH
representativeness sigma}). Under \texttt{network.iono\_supply\_formal\_sigma}
(default OFF, byte-identical) it is root-sum-squared with the SH
coefficients' \emph{formal} $\sigma$, propagated through the same weighted
basis, and the supply is skipped for a satellite that has no estimated
satellite code bias (OSB) in the current network solution.

At the user the rows are not independent. They are folded as a
per-constellation common-mode block with the noise structure of
\dcref{eq:atmos-gim-R}{the GIM iono-weighting covariance}: rows of one
constellation share the variance \texttt{common\_mode\_sigma\_m}$^2$ in
addition to their own diagonal variance. Frozen BSSD
pivots, and any slant state with zero variance, are excluded for the reason
given there: the common-mode block absorbs the datum offset of the
\emph{remaining} rows, so the datum satellite must not appear as a
measurement.

A standardised STEC product on the stream remains open: the in-process
channel is a closed-loop store hand-off, not an IGS-SSR or RTCM message
that a user of a third-party stream can decode.

\subsection{The user of the network's corrections}\label{sec:ssr-netuser}

The network's products reach a user engine in three ways: the in-process
closed loop, in which the network's emit seam fills an owned store each
epoch; the offline replay of a recorded product file
(\texttt{network.closed\_loop.replay\_products\_path}); and, in a real-time
run, a tail of a growing product file
(\texttt{realtime.network\_products\_path}). The first two are described in
\dcref{sec:net-closedloop}{the closed-loop section of the Network mode
document}. Each gives the engine an owned \texttt{SsrCorrectionStore}
(\texttt{RealTimeEngine::setSsrStore}) in place of the process-wide one.

The configuration of the user is derived from the configuration of the network
by \texttt{GnssConfig::\allowbreak deriveClosedLoopConsumerConfig}. The
standalone profile \texttt{processing.\allowbreak network\_consumer} applies
the same couplings (\texttt{applyNetwork\allowbreak ConsumerCouplings}). The
derived configuration differs from that of the network in these respects.
\begin{itemize}[leftmargin=*, itemsep=0pt]
\item The processing mode is the single-receiver PPP-RTK user. Estimation of
      satellite clocks, satellite phase biases, satellite code biases and the
      network ionosphere is switched off, because the corrections supply them.
\item Residual satellite code biases are estimated only when the user does not
      take the network clock (\texttt{estimate\_sat\_biases} is the negation of
      clock consumption). When no satellite-bias state is estimated, the variance
      of a code row on a signal that neither a \texttt{.BIA} entry nor the
      stream covers is multiplied by the ordinary user factor $100$
      (\texttt{no\_bia\_code\_downweight}), in place of the weight the network
      gives it.
\item The undifferenced user applies no observation differencing. The phase
      biases are applied to the raw float \eqref{eq:ssr-phase-bias}, and AR
      forms between-satellite differences at fix time, which cancels the
      receiver phase bias. AR is direct per-band MLAMBDA without the
      Melbourne-W\"ubbena wide lane (\texttt{ar\_wl\_only} false, after which
      the driver sets \texttt{ar\_ignore\_mw\_wl}), and starts only once the
      position $\sigma$ is below $0.2$\,m. The reference configurations run
      this user float (\texttt{network.closed\_loop.ar\_mode} 0).
\item The bias staleness gate is armed under consumption unless the
      configuration sets the key explicitly.
\item Settings that act on the network's own solution (the canonical-signal
      rule, the extra-band code weight, the second code and phase channels) are
      pinned to neutral values, so that a case moves the network and leaves its
      user unchanged.
\end{itemize}
The user's answers to the network's declarations are the phase-bias reactions
of \Cref{sec:ssr-pbstep} and the datum-break reactions of
\dcref{sec:rt-epoch}{the per-epoch update of the Real-time engine document}.

\subsection{Provenance and code pointers}\label{sec:ssr-plumbing}

Both transports map onto the same internal correction structures
(\texttt{rtcm/RtcmTypes.hpp}) through a single decoder
(\texttt{rtcm/RtcmSsrDecoder.cpp}). Legacy RTCM uses per-constellation
six-message blocks (orbit, clock, code bias, combined orbit+clock, URA,
high-rate clock) and a separate phase-bias block (1265--1270).
IGS-SSR~\texttt{4076} repackages the same payloads under a
message indicator (IM). Its \emph{tens} place ($\lfloor\mathrm{IM}/10
\rfloor$) selects the constellation (GPS\,$2x$, GLONASS\,$4x$,
Galileo\,$6x$, QZSS\,$8x$, BeiDou\,$10x$, SBAS\,$12x$), and its
\emph{ones} place ($\mathrm{IM}\bmod 10$) selects the subtype
($1$=orbit, $2$=clock, $3$=combined, $4$=high-rate, $5$=code bias,
$6$=phase bias, $7$=URA). This subtype ordering differs from the legacy
SSR block ordering, so the legacy mapping is not reused
(\texttt{RtcmSsrDecoder.cpp}). A \texttt{4076} message uses one set of field
widths for every constellation (6-bit satellite identifier, 8-bit IOD, 20-bit
GPS-aligned time of week), where the legacy blocks use per-constellation
widths. Applying the legacy Galileo or BeiDou width to a \texttt{4076} message
would over-read the IOD by two bits and shift every later field. RTCM counts
SBAS satellites from $1$ at PRN~120, and the decoders name identifier $n$ as
S$(n+19)$, so PRN~120 is S20 as in RINEX and SBF
(\texttt{rtcm/RtcmSatId.hpp}). VTEC (IM\,$201$) and
NavIC (IM\,$141$--$147$) are rejected by the message-indicator
classifier before the header is parsed, so they return
\texttt{ok = false} with an empty header. For this reason the ionospheric
term receives no SSR contribution from the stream (the
\texttt{SsrSlantIono} store channel of the ionosphere subsection is
unaffected).

The correction definitions
\Cref{eq:ssr-orbit,eq:ssr-clock,eq:ssr-iode-gate} are implemented in the
static, deterministic \texttt{applyOrbitToPosition} / \texttt{applyClockToBias} /
\texttt{iodeMatches} members of \texttt{products/SsrCorrectionStore}, the
interface this document describes. The stateful streaming layer
(the mutex-guarded per-satellite cache with freshness gating on receive
time, the \texttt{RealTimeEngine::feedRtcmFrame} push loop, and the
obs$\leftrightarrow$SSR synchronisation that bounds how far the obs
stream may run ahead of lagging corrections) is described in
\dcref{sec:rt}{the Real-time engine document}. It is kept separate so that
the correction mathematics can be analysed and tested in isolation. The
estimator-side warm-up gate
\texttt{estimation.defer\_until\_corrections\_ready.enabled}
(\texttt{config/GnssConfigLoad.cpp}; default OFF, real-time only) defers
estimation until a quorum of satellites is centimetre-grade from broadcast
plus SSR (by default at least 5 in total and 4 per constellation). The
filter then seeds code and phase together, instead of cold-starting on
metre-grade broadcast clocks and fixing a contaminated
ionosphere/bias/ambiguity partition.

The store also carries a per-epoch \emph{heartbeat} of the network
(\texttt{NetworkProducerStatus}). It lets a user distinguish ``the
network withheld this epoch'' from ``the network stopped'': a withheld
epoch still writes its status line, although it publishes no corrections.
The token set of the product file is forward-compatible: an unrecognised
token from a newer network maps to \texttt{State::Unknown} and is reported,
never fatal (\texttt{NetworkProducerStatus::parseState} / \texttt{stateName},
declared in the store's header). The token \texttt{withheld\_sh} is retained
so that archived product files parse; the spherical-harmonic health gate that
wrote it was deleted from the network, and no current product contains it.

\ifSubfilesClassLoaded{\printbibliography}{}
\end{document}
